\section{Explicit kings: when is the frozen tr\'ebuchet exact?}\label{sec:kings}

The convention treats a king group as $0$; \citet{Maeser2002} already warned that this is heuristic, and \citet{Siegel2009cycles} studied loopy chess positions. We put the kings back with an exact king-and-pawn solver,
\file{kpsolve}. It plays real chess with kings, pawns, en passant, all four promotions and stalemate, and no
fifty-move rule. An independent reference solver agrees with it on all $1{,}418{,}480$ positions of $18$
small-board runs. On every comparable position it is consistent with the Lichess tablebase
(\status{checked-2}).

\subsection{Validity conditions}
\begin{theorem}[\status{proved}; reviewed]\label{thm:validity}
Let a position consist of pawn components and a king group. Its chess outcome equals the normal-play outcome
of the pawn components if all of the following hold:
\begin{enumerate}[label=(V\arabic*)]
\item the kings never interact with the pawn components, and nobody is ever stalemated;
\item every \emph{decisive event} (the appearance of a passed pawn) really wins;
\item every move inside the king group is decisively punished.
\end{enumerate}
The step from there to the $\pm\BIG$ arithmetic needs one more condition:
\begin{enumerate}[label=(V\arabic*),resume]
\item no two components carry passed-pawn threats of opposite sign, and $\BIG$ exceeds every tempo count.
\end{enumerate}
\end{theorem}
(V3) has a local certificate: the punishing side promotes within $T$ moves even if the defender may pass,
while the defender cannot promote within $T$ moves.
\begin{itemize}
\item With $T=1$ the certificate holds in every context.
\item With $T\ge2$ it fails in some context.
\end{itemize}
Real tr\'ebuchets have $T\approx5$--$7$, so their ``zero'' is \emph{context-dependent}: it is valid only beside
slow or locked components.

\subsection{Spare tempo moves are loopy}
A king that can shuttle forever without loss is not $0$ but a loopy game:
\begin{itemize}
\item $\textsf{on}=\game{\textsf{on}}{}$ if only White can shuttle;
\item $\textsf{off}=\game{}{\textsf{off}}$ if only Black can;
\item $\textsf{dud}$ if both can.
\end{itemize}
Since $\textsf{on}$ exceeds every number, $\BIG$ included, the $\BIG$ arithmetic breaks down completely there.
\begin{theorem}[\status{proved}]\label{thm:loopy}
\begin{itemize}
\item If only White can pass, Black wins iff Black can force a decisive event against a passing White;
  otherwise White wins.
\item If both sides can pass, a side wins iff it can force a decisive event against a passing opponent;
  otherwise the game is drawn.
\end{itemize}
\end{theorem}

\subsection{The paper's king groups in real chess}
\begin{table}[ht]
\centering\small
\begin{tabularx}{\textwidth}{@{}l l l l X@{}}
\toprule
Group & Board & White to move & Black to move & Verdict \\
\midrule
5.1: Kh5 Pa4 / Kb4 Pa5 & cylinder & Black wins & White wins & real mutual zugzwang (\status{checked-2}) \\
\cref{fig:candidate}: Kf5 Pg4 / Kh4 Pg5 & flat & Black wins & White wins & real mutual zugzwang (Lichess) \\
7.2: Ke6 Pd7 / Kd8 Pe7 & flat & \textbf{draw} & White wins & \textbf{not} a zugzwang (Lichess) \\
\bottomrule
\end{tabularx}
\caption{King groups standing alone. The convention reads each as $0$.}
\label{tab:kings}
\end{table}
In 7.2, White to move draws with Kd5, Ke5 or Kf5: the d7 pawn falls, but the white king catches the e-pawn.
Only Kf7 loses. The audit's statement that 7.2 is a mutual zugzwang is therefore true only inside the
convention.

\subsection{How often does the convention fail?}
In sweeps over every small pawn component beside a tr\'ebuchet, on $6\times6$ and $7\times6$ cylinders and a
$5\times6$ flat board, $159$ of $500$ testable components get the wrong outcome (\status{computed-1}).
\begin{itemize}
\item In $107$ the kings alone are the cause: with frozen kings the solver reproduces \file{cylcgt} exactly.
\item The other $52$, all on $7\times6$, are failures of the pawn convention itself: $18$ through en passant,
  and $34$ because a passed pawn does not actually win its race.
\end{itemize}
Small boards are a stress test, since the kings are never far from the pawns. On $8\times8$ the failure rate is
lower but not zero (\cref{sec:star2}).
